\section{Theoretical Foundation: Four Core Principles}

\subsection{Principle 1: Capability Abstraction Principle}

\paragraph{Principle:}
\emph{An AI agent should depend on capabilities rather than concrete providers or APIs.}

\paragraph{Formal Definition:}
A capability $C$ is defined as:
\begin{equation}
C = (I, O, \Sigma, Q, P)
\label{eq:capability-tuple}
\end{equation}
where:
\begin{itemize}[leftmargin=*,nosep]
  \item $I$: input schema
  \item $O$: output schema
  \item $\Sigma$: execution semantics
  \item $Q$: quality telemetry schema (e.g., latency, retry\_count, reliability signals)
  \item $P$: pricing and billing schema (e.g., unit price, billing mode, usage counters)
\end{itemize}

A Provider $P$ implements Capability $C$. Agents invoke capabilities through
\begin{equation}
\mathrm{Invoke}(C, \mathrm{Input}) \rightarrow \mathrm{Output}.
\label{eq:invoke}
\end{equation}

Rather than invoking provider-specific APIs directly.

\paragraph{Significance:}
This principle achieves agent-provider decoupling. Agents express what they need (capabilities) rather than how to obtain it (specific APIs). This abstraction enables:
\begin{itemize}[leftmargin=*,nosep]
  \item Provider substitution without agent modification
  \item Multi-provider capability implementations
  \item Standardized capability discovery and invocation
\end{itemize}

\begin{table}[t]
\caption{Comparison between capability and related abstractions.}
\label{tab:capability-comparison}
\centering
\small
\begin{tabularx}{\linewidth}{@{} l X X @{}}
\toprule
Concept & Essence & Role \\
\midrule
Function & Computation & Algorithmic operation \\
API & Remote interface & Access mechanism \\
Tool & Agent-callable function & Invocation interface \\
\textbf{Capability} & \textbf{Executable ability} & \textbf{Real-world action} \\
\bottomrule
\end{tabularx}
\end{table}

Capabilities represent abilities (what can be done), while APIs represent access mechanisms (how to invoke). This distinction is fundamental to the AI Native Architecture.

In this paper's scope, the execution runtime does not optimize $Q$ or $P$; it only measures, normalizes, and forwards these fields for orchestration-layer policy decisions.

\subsection{Principle 2: Cognition--Action Separation Principle}

\paragraph{Principle:}
\emph{AI cognition and capability execution should be separated by an action layer.}

\paragraph{Architecture:}
\begin{center}
\begin{tabular}{c}
Cognition Layer (reasoning, planning, memory) \\
$\downarrow$ \\
Action Layer (routing, execution) \\
$\downarrow$ \\
Capability Providers
\end{tabular}
\end{center}

\paragraph{Rationale:}
Without this separation, agents must directly call APIs, tools, and services, leading to:
\begin{itemize}[leftmargin=*,nosep]
  \item Tight coupling to specific providers
  \item Low portability across different execution environments
  \item Ecosystem fragmentation
\end{itemize}

The action layer mediates between cognitive processes (deciding what to do) and execution processes (performing the action), enabling agents to focus on reasoning while the execution layer handles the mechanics of capability invocation.

\subsection{Principle 3: Capability Routing Principle}

\paragraph{Principle:}
\emph{Capability execution should be dynamically routed based on context.}

\paragraph{Routing Function:}
\begin{equation}
\mathrm{Route}(C, \mathrm{Context}) \rightarrow \mathrm{Provider}
\label{eq:routing}
\end{equation}

\paragraph{Context Parameters:}
\begin{itemize}[leftmargin=*,nosep]
  \item \textbf{Latency:} Response time requirements
  \item \textbf{Cost:} Budget constraints
  \item \textbf{Policy:} Organizational or regulatory policies
  \item \textbf{Availability:} Provider uptime and reliability
  \item \textbf{Security:} Data handling and privacy requirements
\end{itemize}

\paragraph{Example:}
For a capability \texttt{translate}, available providers might include:
\begin{itemize}[leftmargin=*,nosep]
  \item Provider A: Lower cost, higher latency
  \item Provider B: Higher cost, lower latency
\end{itemize}

The routing system dynamically selects based on context. If the user priority is speed, Provider B is selected. If cost optimization is paramount, Provider A is chosen.

\paragraph{Analogy:}
This mechanism is similar to DNS routing or service mesh routing~\cite{consul_service_mesh}, but the objects being routed are capabilities rather than network requests.

\subsection{Principle 4: Capability Graph Execution}

\paragraph{Principle:}
\emph{Complex AI tasks can be modeled as execution over a capability graph.}

\paragraph{Formal Definition:}
A Capability Graph $G$ is defined as:
\begin{equation}
G = (\mathrm{Capabilities}, \mathrm{Dependencies})
\label{eq:capability-graph}
\end{equation}

For execution safety and deterministic ordering, we constrain $G$ to a directed acyclic graph (DAG). Cycles are disallowed.

\paragraph{Example:}
The task ``Write and publish a blog post'' corresponds to the capability graph shown in \cref{fig:capability-graph}.

\begin{figure}[t]
\centering
\begin{tikzpicture}[
  node distance=1.2cm and 0.8cm,
  capnode/.style={draw, rounded corners=2pt, minimum width=2.0cm, minimum height=0.8cm, align=center, fill=blue!6},
  exec/.style={draw, rounded corners=2pt, minimum width=1.2cm, minimum height=0.6cm, align=center, fill=green!8},
  >=Latex
]
\node[capnode] (search) {\texttt{search}};
\node[capnode, right=of search] (sum) {\texttt{summarize}};
\node[capnode, right=of sum] (gen) {\texttt{generate\_text}};
\node[capnode, right=of gen] (trans) {\texttt{translate}};
\node[capnode, right=of trans] (pub) {\texttt{publish}};

\node[exec, below=of search] (e1) {exec};
\node[exec, below=of sum] (e2) {exec};
\node[exec, below=of gen] (e3) {exec};
\node[exec, below=of trans] (e4) {exec};
\node[exec, below=of pub] (e5) {exec};

\draw[->] (search) -- (sum);
\draw[->] (sum) -- (gen);
\draw[->] (gen) -- (trans);
\draw[->] (trans) -- (pub);

\draw[->] (search) -- (e1);
\draw[->] (sum) -- (e2);
\draw[->] (gen) -- (e3);
\draw[->] (trans) -- (e4);
\draw[->] (pub) -- (e5);
\end{tikzpicture}
\caption{Capability graph for a multi-step publishing task.}
\label{fig:capability-graph}
\end{figure}

\paragraph{Execution Process:}
\begin{enumerate}[leftmargin=1.5em]
    \item Agent reasons about the task and constructs the capability graph
    \item Orchestration layer determines execution order and provider selection
    \item Execution runtime invokes each capability in sequence
    \item Results flow through the graph, with each capability receiving inputs from predecessors
\end{enumerate}

This model enables complex tasks to be decomposed into composable capability sequences, with the execution runtime handling the mechanics of each capability invocation.
